\section*{Bab \ref{ch:b3:lab-techniques-3} --- Teknik Laboratorium III: Praktik Penelitian}

\begin{solution}{exo:b3:lab-techniques-3:1}
$(0.50/1013)^3 = \num{1.2e-10}$. Kebocoran, gas yang teradsorpsi pada kaca, serta pelepasan gas dari gemuk dan septum
menyisakan jauh lebih banyak daripada itu; peralatan kaca yang kering dan panas serta sambungan yang baik lebih penting daripada
siklus tambahan.
\end{solution}

\begin{solution}{exo:b3:lab-techniques-3:2}
$\qty{20.0}{mmol}/\qty{1.48}{mol/L} = \qty{13.5}{mL}$.
\end{solution}

\begin{solution}{exo:b3:lab-techniques-3:3}
$120 + 10 \times 1.00782503 + 55.93493633 - 0.00054858 = \qty{186.01264}{u}$. Galat: $10^6 \times (186.0129 - 186.01264)/186.01264 = +1.4$ ppm,
jauh di dalam batas 5 ppm.
\end{solution}

\begin{solution}{exo:b3:lab-techniques-3:4}
Primer: artikel jurnal, paten, tesis. Sekunder: ulasan,
buku pegangan tetapan. Tersier: buku teks.
\end{solution}

\begin{solution}{exo:b3:lab-techniques-3:5}
$M = \qty{194.0}{g/mol}$: C $96.0/194.0 = 49.48$\,\%, H $10.0/194.0 = 5.15$\,\%, N $56.0/194.0 = 28.87$\,\%. Selisihnya 0.17, 0.07, dan 0.16
poin: di dalam batas 0.4, sehingga analisis itu mendukung rumusnya.
\end{solution}

\begin{solution}{exo:b3:lab-techniques-3:6}
$\qty{170.0}{mg}/\qty{212.0}{g/mol} = \qty{0.802}{mmol}$; $c = \qty{0.802}{mmol}/\qty{0.54}{mL} = \qty{1.5}{mol/L}$ (dua angka penting, yang ditentukan oleh
pembacaan volume).
\end{solution}

\begin{solution}{exo:b3:lab-techniques-3:7}
Rata-ratanya 73.0 dan 77.0\,\%; simpangan bakunya sama-sama $\sqrt{10/4} = 1.58$; $s_p = 1.58$. $t = 4.0/(1.58\sqrt{2/5}) = 4.0$, lebih besar daripada
2.31: perbaikan itu signifikan pada tingkat 95\,\%.
\end{solution}

\begin{solution}{exo:b3:lab-techniques-3:8}
$F = (0.80/0.30)^2 = 7.1 > 5.05$: metode pertama secara signifikan kurang presisi.
\end{solution}

\begin{solution}{exo:b3:lab-techniques-3:9}
Belilah botol-botol kecil yang berinhibitor, beri tanggal pada setiap botol saat diterima dan saat dibuka, simpanlah
tertutup, di tempat gelap dan sejuk; ujilah kadar peroksidanya sebelum distilasi atau penguapan apa pun
dan pada selang waktu tetap setelah dibuka; buanglah botol yang telah lewat tanggalnya, atau yang
positif, melalui rute limbah berbahaya tanpa memekatkannya.
\end{solution}

\begin{solution}{exo:b3:lab-techniques-3:10}
Bahaya: hidrogen yang dilepaskan (kebakaran, tekanan), kereaktifan NaH dengan air,
dan penguraian eksotermik DMF bersama natrium hidrida yang sudah dikenal, yang dapat lepas kendali
jika dipanaskan; DMF juga beracun bagi reproduksi. Pengendalian: gantilah pelarutnya
(THF atau 2-metiltetrahidrofuran) atau basanya; jika tetap dipakai, jalankan dalam skala kecil lebih dahulu, pada suhu
rendah, dengan menambahkan NaH sedikit demi sedikit sambil memantau suhu, di dalam lemari asam
di balik perisai, dengan saluran buang melalui gelembung, dan dengan tangas pendingin yang siap; prosedur
tertulis, orang kedua yang hadir, dan alat pelindung diri sebagai upaya terakhir.
\end{solution}

\begin{solution}{exo:b3:lab-techniques-3:11}
Kontribusi relatif: 0.50, 0.50, $0.010/20.0 = 0.05$, $0.010/15.0 = 0.067$, dan 0.10\,\%; gabungannya
$\sqrt{0.25 + 0.25 + 0.0025 + 0.0044 + 0.010} = 0.72$\,\%. Kedua integral itu dominan: rasio sinyal terhadap derau yang lebih baik serta
koreksi garis dasar dan fase yang cermat lebih penting daripada neraca.
\end{solution}

\begin{solution}{exo:b3:lab-techniques-3:12}
$Q = (89.0 - 82.0)/(89.0 - 80.6) = 0.83 > 0.71$: uji itu menandainya. Sebelum menolaknya, periksalah buku catatan
untuk mencari penyebabnya (buret yang salah dibaca, gelembung udara, batch titran yang berbeda); jika penyebab itu
ditemukan, tolaklah nilai itu dan sebutkan alasannya; jika tidak, pertahankan nilai itu, laporkan ujinya, atau ulangi
pengukurannya.
\end{solution}

\begin{solution}{pb:b3:lab-techniques-3:1}
\textbf{1.} $\ce{FeCl2 + 2NaC5H5 -> Fe(C5H5)2 + 2NaCl}$.

\textbf{2.} Natrium siklopentadienida dirusak oleh oksigen dan air, dan besi(II)
klorida teroksidasi menjadi besi(III) di udara lembap.

\textbf{3.} $\qty{5.00}{g}/\qty{126.8}{g/mol} = \qty{0.0394}{mol}$; $\times \qty{185.8}{g/mol} = \qty{7.33}{g}$; rendemen $5.86/7.33 = 80$\,\%.

\textbf{4.} $(0.10/1000)^3 = 10^{-12}$ (secara teoretis).

\textbf{5.} Dengan memanaskan dimernya, yang mengalami reaksi retro-Diels--Alder, dan
mendistilasi monomernya yang atsiri; pada suhu kamar monomer itu berdimerisasi
lagi melalui reaksi Diels--Alder, sehingga monomer itu disimpan dingin dan segera dipakai.

\textbf{6.} Setiap porsi melepaskan hidrogen dan kalor: menambahkannya perlahan-lahan menjaga aliran gas
dan suhu tetap terkendali.

\textbf{7.} \qty{186.0126}{u}; galat $+1.4$ ppm.

\textbf{8.} C $120.0/185.8 = 64.58$\,\%, H $10.0/185.8 = 5.38$\,\%; yang ditemukan 64.41 dan 5.47: selisihnya $-0.17$ dan $+0.09$, di dalam batas 0.4.

\textbf{9.} Kesepuluh hidrogen setara karena simetri, dan cincin-cincinnya berputar dengan cepat.

\textbf{10.} Satu sinyal, karena kesepuluh karbon setara.

\textbf{11.} Struktur sinar-X kristal tunggal.

\textbf{12.} Tidak: sebagai kompleks 18 elektron (\cref{ch:b3:organometallic-bonding}), ferosena stabil di udara, tidak seperti
prekursor-prekursornya.

\textbf{13.} Sinyal-sinyalnya tidak bertumpang tindih dengan sinyal ferosena, zat itu berupa padatan murni yang tidak atsiri dan
tidak higroskopis sehingga dapat ditimbang dengan teliti, dan zat itu inert terhadap
sampel.

\textbf{14.} Relaksasi penuh setiap sinyal di antara pemindaian (jeda beberapa kali $T_1$),
eksitasi yang seragam, rasio sinyal terhadap derau yang cukup, serta koreksi fase dan garis dasar
yang cermat.

\textbf{15.} Ferosena $\ce{C10H10Fe}$ \qty{185.8}{g/mol}; 1,3,5-trimetoksibenzena $\ce{C9H12O3}$ \qty{168.0}{g/mol}.

\textbf{16.} $P_x = 3.942 \times (3/10) \times (185.8/168.0) \times (15.0/20.0) \times 0.999 = 0.980$: 98.0\,\%.

\textbf{17.} Relatif $\sqrt{0.50^2 + 0.50^2 + 0.05^2 + 0.067^2 + 0.10^2} = 0.72$\,\%; absolut $0.72\,\% \times 98.0 = 0.7$ poin persentase.

\textbf{18.} Ya, tetapi lemah: pengotor 2\,\% dengan komposisi serupa mengubah C dan H kurang
dari yang dapat dideteksi analisis itu; qNMR adalah pengukuran kemurnian yang lebih informatif.
\textbf{19.} Natrium hidrida (hidrogen dengan air, kebakaran), hidrogen yang dilepaskan, THF dan
siklopentadiena (sangat mudah menyala; THF \omterm{def:b3:lab-techniques-3:peroxide}{pembentuk peroksida}), perengkahan panas
dimer, besi(II) klorida (iritan), dan ferosena (padatan mudah menyala, berbahaya).

\textbf{20.} Natrium hidrida: ditimbang sebagai dispersi dalam minyak di bawah gas inert, ditambahkan sedikit
demi sedikit ke dalam larutan yang diaduk dan didinginkan, dengan gas dibuang melalui gelembung, tanpa air
di dekatnya. THF: kering dan bebas peroksida dari kolom, dipakai di dalam lemari asam di bawah nitrogen.

\textbf{21.} Di bawah gas inert dan dengan pendinginan, dengan penambahan isopropanol secara perlahan, lalu
etanol, lalu air, dengan menunggu setiap kali sampai pelepasan gas berhenti.

\textbf{22.} Gelombang satu-elektron yang reversibel, $\ce{Fe(C5H5)2+}/\ce{Fe(C5H5)2}$, dengan puncak-puncak yang terpisah sekitar 59 mV pada \qty{25}{\celsius}
(\cref{ch:b3:electrode-kinetics}); pasangan itu cepat, stabil, dan hampir tidak peka terhadap
pelarut, sehingga menjadi acuan internal yang praktis.

\textbf{23.} Buku catatan: tanggal, skema, jumlah dan sumber, \omterm{def:b3:lab-techniques-3:assessment}{penilaian risiko}, apa yang
dilakukan dan dilihat, rendemen, dan nama berkas \omterm{def:b3:lab-techniques-3:notebook}{data mentah}. \omterm{def:b3:lab-techniques-3:literature}{Informasi
pendukung}: prosedur lengkap dan semua data karakterisasi beserta salinan
spektrumnya.

\textbf{24.} Agar orang lain, dan para penulis bertahun-tahun kemudian, dapat membuka kembali dan mengolah ulang data itu
dengan perangkat lunak apa pun yang mereka miliki: mudah ditemukan, dapat diakses, dapat saling dioperasikan, dan dapat dipakai ulang.

\textbf{25.} Kemurnian qNMR sampel ferosena adalah $98.0 \pm 0.7$\,\% (ketidakpastian baku).
\end{solution}
